\documentclass[14pt]{beamer}
\usepackage[T2A]{fontenc}
\usepackage[utf8]{inputenc}
\usepackage[english,russian]{babel}
\usepackage{graphicx}

\usepackage{fontspec}
\setmainfont{Times New Roman}
\setsansfont{Arial}
\setmonofont{Courier New}

\usetheme{Madrid}
\usecolortheme{default}

\title{Лабораторная работа №7}
\subtitle{Создание презентаций в LaTeX с помощью Beamer}
\author{Сунь Маосин}
\institute{Российский университет дружбы народов}
\date{2026}

\begin{document}

% Титульный слайд
\begin{frame}
  \titlepage
\end{frame}

% Содержание
\begin{frame}
  \frametitle{Содержание}
  \tableofcontents
\end{frame}

% Раздел 1
\section{Введение}

\begin{frame}{Что такое Beamer?}
  \begin{block}{Определение}
    Beamer — это класс LaTeX для создания презентаций.
    Он позволяет создавать профессиональные слайды с математикой, графикой и анимацией.
  \end{block}
  
  \begin{exampleblock}{Преимущества}
    \begin{itemize}
      \item Красивое оформление
      \item Поддержка математических формул
      \item Лёгкость в использовании
    \end{itemize}
  \end{exampleblock}
\end{frame}

% Раздел 2
\section{Структура презентации}

\begin{frame}{Основные элементы}
  \begin{columns}
    \column{0.5\textwidth}
      \begin{block}{Заголовки}
        Каждый слайд начинается с 
        \texttt{\textbackslash begin\{frame\}} 
        и заканчивается 
        \texttt{\textbackslash end\{frame\}}.
      \end{block}
    
    \column{0.5\textwidth}
      \begin{block}{Секции}
        \texttt{\textbackslash section} создаёт разделы,
        которые автоматически появляются в содержании.
      \end{block}
  \end{columns}
\end{frame}

% Раздел 3
\section{Списки и анимация}

\begin{frame}{Появление элементов}
  \begin{itemize}
    \item Первый пункт появляется сразу
    \pause
    \item Второй пункт появляется после клика
    \pause
    \item Третий пункт появляется ещё позже
    \pause
    \item И так далее...
  \end{itemize}
\end{frame}

\begin{frame}{Управление появлением}
  \begin{enumerate}[<+->]
    \item Этот пункт появляется первым
    \item Этот — вторым
    \item Этот — третьим
  \end{enumerate}
\end{frame}

% Раздел 4
\section{Математика в презентациях}

\begin{frame}{Формулы}
  \begin{block}{Знаменитая формула Эйнштейна}
    \[
      E = mc^2
    \]
  \end{block}
  
  \begin{block}{Интеграл}
    \[
      \int_{0}^{\infty} e^{-x^2} dx = \frac{\sqrt{\pi}}{2}
    \]
  \end{block}
\end{frame}

% Раздел 5
\section{Таблицы и графики}

\begin{frame}{Простая таблица}
  \begin{table}
    \centering
    \begin{tabular}{|l|c|r|}
      \hline
      \textbf{Город} & \textbf{Население} & \textbf{Страна} \\
      \hline
      Москва & 12.5 млн & Россия \\
      Пекин & 21.5 млн & Китай \\
      Вашингтон & 0.7 млн & США \\
      \hline
    \end{tabular}
    \caption{Население городов}
  \end{table}
\end{frame}

% Раздел 6
\section{Разные темы}

\begin{frame}{Доступные темы}
  В Beamer есть много готовых тем:
  \begin{columns}[t]
    \column{0.3\textwidth}
      \begin{itemize}
        \item Madrid
        \item Copenhagen
        \item Warsaw
      \end{itemize}
    \column{0.3\textwidth}
      \begin{itemize}
        \item Berlin
        \item Dresden
        \item Frankfurt
      \end{itemize}
    \column{0.3\textwidth}
      \begin{itemize}
        \item Malmoe
        \item Montpellier
        \item Berkeley
      \end{itemize}
  \end{columns}
\end{frame}

% Раздел 7
\section{Заключение}

\begin{frame}{Вывод}
  \begin{block}{Что мы узнали}
    \begin{itemize}
      \item Как создавать слайды в Beamer
      \item Как использовать разные темы
      \item Как добавлять списки и анимацию
      \item Как вставлять математические формулы
      \item Как делать таблицы и колонки
    \end{itemize}
  \end{block}
  
  \begin{alertblock}{Важно!}
    Beamer — отличный инструмент для научных презентаций.
    Он экономит время и даёт профессиональный результат.
  \end{alertblock}
\end{frame}

% Последний слайд
\begin{frame}
  \centering
  \Huge Спасибо за внимание!
  
  \vspace{1cm}
  \large Вопросы?
\end{frame}

\end{document}